\section{ECNetBench: Benchmark Design}
\label{sec:bench}

\subsection{Design goals}
ECNetBench targets enterprise cognitive networking research with five goals: (G1)~AOS-CX-style multi-layer configuration realism; (G2)~multi-modal telemetry (counters, resources, syslog/alerts, flow aggregates); (G3)~causal incident chains linking failure, detection, recovery, and impact; (G4)~leakage-safe supervised labels; (G5)~multi-seed robustness under frozen instances. Table~\ref{tab:contrib} contrasts ECNetBench with common public resource classes for cognitive NetOps research.
\begin{table}[!t]
\centering
\caption{Qualitative comparison of ECNetBench against common public resource classes for enterprise cognitive NetOps research.}
\label{tab:contrib}
\renewcommand{\arraystretch}{1.15}
\setlength{\tabcolsep}{3.5pt}
{\scriptsize\begin{tabular}{p{1.85cm}cccc}
\toprule
Resource class & Config. & Telem. & RCA/heal & Leak.-safe multi-seed \\
\midrule
Topology archives & Partial & No & No & Rare \\
Security flow sets & Low & Flows & Limited & Varies \\
Clos/telemetry RCA & Fabric & Yes & RCA & Varies \\
\textbf{ECNetBench} & \textbf{Yes} & \textbf{Yes} & \textbf{Yes} & \textbf{Yes} \\
\bottomrule
\end{tabular}}
\end{table}


\subsection{Enterprise assumptions and topology}
Instances model a compact enterprise fabric (19 devices, 31 links in the evaluated family) spanning campus/DC roles with VLAN membership, ACL/QoS bindings, routing (e.g., BGP/OSPF-related tables), LAG/VSX-like redundancy concepts, and overlay-related objects where applicable. Topology snapshots and typed graph edges support Digital Twin construction. Generation constraints and schema documents accompany the release.

\subsection{Telemetry and causal incident modeling}
Synthetic generators emit interface counters, device resources, environmental/power samples, syslog/alert streams, and service KPIs. Failure incidents are drawn from a taxonomy covering link, device, control-plane, policy, and service-impacting families. Each incident carries detection, recovery-action, and impact fields so that labels for anomaly windows, failure horizons, RCA categories, impact, and healing actions remain consistent with a causal narrative. Fig.~\ref{fig:benchflow} summarizes the generation and validation workflow from constrained synthesis to audited SQLite instances and supervised labels.

\begin{figure}[!t]
\centering
\begin{tikzpicture}[
  box/.style={draw,align=center,font=\footnotesize\sffamily,inner sep=3.5pt,minimum height=0.78cm,minimum width=2.05cm},
  arr/.style={-{Latex},thick}
]
\node[box,fill=black!6] (gen) {Generator\\+ constraints};
\node[box,right=0.5cm of gen,fill=black!4] (db) {SQLite\\instance};
\node[box,right=0.5cm of db,fill=black!4] (lab) {Labels\\T1--T6};
\node[box,below=0.48cm of db,fill=black!4] (aud) {Realism +\\leakage audits};
\draw[arr] (gen) -- (db);
\draw[arr] (db) -- (lab);
\draw[arr] (db) -- (aud);
\end{tikzpicture}
\caption{ECNetBench generation and validation workflow: constrained synthesis produces a frozen SQLite instance that is audited for realism/leakage and exported to supervised task labels.}
\label{fig:benchflow}
\end{figure}

\subsection{Leakage-safe labels and splits}
Labels are aligned to time bins. Features at decision time $t_0$ use only information available at or before $t_0$ (historical shifts; no future counters). Evaluation uses a temporal $70/15/15$ train/validation/test split frozen per instance. Forbidden fields (raw future labels, oracle incident IDs in feature matrices) are audited in the companion leakage reports.

\subsection{Instances and seed robustness}
We evaluate six frozen SQLite instances: primary instance v1.1.0-INST and independent seeds 101, 202, 303, 404, and 505 (Table~\ref{tab:seeds}). SHA-256 digests are published in the instance checksum manifest. Seed-robustness and publication-validation reports accompany the package. SQLite is authoritative for evaluation; CSV/Parquet exports are optional derivatives. Schema families are summarized in Table~\ref{tab:schema}.
\begin{table}[!t]
\centering
\caption{Frozen evaluation instances used in this study.}
\label{tab:seeds}
\renewcommand{\arraystretch}{1.15}
{\footnotesize\begin{tabular}{llc}
\toprule
Instance folder & Role & Devices / links \\
\midrule
\texttt{v1} (v1.1.0-INST) & Primary frozen instance & 19 / 31 \\
\texttt{v1.1-seed101} & Independent seed & 19 / 31 \\
\texttt{v1.1-seed202} & Independent seed & 19 / 31 \\
\texttt{v1.1-seed303} & Independent seed & 19 / 31 \\
\texttt{v1.1-seed404} & Independent seed & 19 / 31 \\
\texttt{v1.1-seed505} & Independent seed & 19 / 31 \\
\bottomrule
\end{tabular}}
\end{table}

\begin{table}[!t]
\centering
\caption{Summary of ECNetBench relational schema families.}
\label{tab:schema}
\renewcommand{\arraystretch}{1.15}
{\footnotesize\begin{tabular}{p{2.4cm}p{5.3cm}}
\toprule
Family & Representative contents \\
\midrule
Inventory/topology & Devices, interfaces, topology snapshots, graph nodes/edges \\
Configuration & VLAN/ACL/QoS, routing processes, configuration snapshots/diffs \\
Telemetry & Interface counters, resources, syslog/alerts, flow aggregates \\
Services & Applications, services, SLA and impact bindings \\
Incidents/labels & Failure incidents, recovery actions, supervised label tables \\
\bottomrule
\end{tabular}}
\end{table}

